% ~200 words. Mirrors the final chapter set; numbers from DATA-AUDIT.md.
NebulaStream, a data management system for the Internet of Things, stores operator state in main memory.
When the state of windowed aggregations grows beyond the available memory, the operating system terminates the process, halting the stream processing.
This thesis builds and evaluates an out-of-core state backend for NebulaStream, which spills window state to disk when memory is scarce.
RocksDB serves as the embedded key-value store. Spills are triggered under soft and hard memory thresholds, and state is restored for result emission.

The evaluation comprises about 3{,}100 measured runs on synthetic sensor streams, under hard memory caps of 512 to 2048\,MB, with every follow-up hypothesis pre-registered.
Within a characterized operating window, the backend completes workloads whose state exceeds memory where the in-memory engine always fails (14/24 versus 0/18 at 512\,MB, Fisher $p = 4.1 \times 10^{-5}$), with results verified by checksum against an unconstrained baseline.
At a matched I/O buffer size, for workloads that both the backend and the in-memory engine complete, the throughput cost is 11--16\%.
The evaluation further identifies an engine default that excluded all window state from the buffer pool.
A single configuration change enables the intended pooling; combined with new pool-pressure levers, it recovers most of the previously failing workloads and, at one measured operating point, reaches 98\% of the in-memory engine's throughput at 25\% less anonymous memory.
The choice between pooled and heap allocation also determines how the buffer pool must be sized.
The completion result transfers beyond these synthetic workloads: on a Nexmark bid stream whose state exceeds a 1024\,MB cap, the pooled backend completes where the in-memory engine cannot.
